\chapter{Cantera in practice}
\label{ch:cantera}

Cantera is an open-source C++ library with a Python interface that evaluates everything in
Chapters~\ref{ch:thermo}--\ref{ch:transport} (thermodynamics, kinetics, transport) and solves the model
problems of Chapters~\ref{ch:reactor}, \ref{ch:premixed} and \ref{ch:nonpremixed}. This repository uses
Cantera 3.2.0 on Python 3.14.

\section{Setup check}
\begin{lstlisting}
source .venv/bin/activate
python src/check_setup.py
# Cantera 3.2.0: adiabatic flame T = 2225.5 K
# Laminar flame speed = 0.419 m/s
\end{lstlisting}

\section{The \texttt{Solution} object}
\label{sec:ct-solution}

\begin{figure}[H]
\centering
\begin{tikzpicture}[node distance=5mm and 8mm]
  \node[boxnode=gray,text width=3.4cm] (y) {mechanism file\\\code{h2o2.yaml}};
  \node[boxnode=intu,right=14mm of y,text width=4.6cm,minimum height=3.2cm] (s) {};
  \node[font=\bfseries\small] at ($(s.north)+(0,-0.3)$) {\code{gas = ct.Solution(...)}};
  \node[boxnode=keyc,text width=4.0cm] at ($(s.center)+(0,0.25)$) (th) {\textbf{ThermoPhase}: state, $h$, $s$, $c_p$, \code{equilibrate}};
  \node[boxnode=hot,text width=4.0cm,below=1mm of th] (ki) {\textbf{Kinetics}: $k_f$, $q_i$, $\wdot_k$, HRR};
  \node[boxnode=cold,text width=4.0cm,below=1mm of ki] (tr) {\textbf{Transport}: $\mu$, $\lambda$, $D_{km}$};
  \draw[->,thick] (y) -- (s);
  \node[boxnode=proj,right=14mm of s,text width=3.2cm] (r) {Reactors (0D)\\FreeFlame (1D)\\Counterflow (1D)};
  \draw[->,thick] (s) -- node[above,lbl] {used by} (r);
\end{tikzpicture}
\caption{A \code{Solution} bundles three managers that all read the same thermodynamic state.
Reactors and flames call them at every point and time step.}
\end{figure}

\begin{lstlisting}
import cantera as ct
gas = ct.Solution("h2o2.yaml")                 # 10 species, 29 reactions
gas.TPX = 300.0, ct.one_atm, "H2:2, O2:1, N2:3.76"   # set state (T, p, composition)

gas.T, gas.P, gas.density                      # state
gas.X, gas.Y, gas.species_names                # composition
gas.enthalpy_mass, gas.cp_mass, gas.entropy_mass       # J/kg, J/kg/K
gas.net_production_rates, gas.heat_release_rate        # kmol/m3/s, W/m3
gas.viscosity, gas.thermal_conductivity, gas.mix_diff_coeffs_mass
\end{lstlisting}

Setters come in pairs: \code{TP}, \code{TPX}, \code{TPY}, \code{HP}, \code{HPY}, \code{UV}, \code{SP}, \dots
\emph{All} properties update consistently with the state.

\begin{warning}[Python 3.14 / NumPy 2.5 quirk]
In this environment \code{gas["OH"].X} raises
\code{ValueError: cannot resize an array that may be referenced}. Use index access instead:
\begin{lstlisting}
iOH = gas.species_index("OH")
x_OH = gas.X[iOH]
\end{lstlisting}
The project page uses \code{gas["OH"].X[0]}; translate it this way.
\end{warning}

\section{Anatomy of a YAML mechanism}
\label{sec:ct-yaml}

Excerpts from \code{h2o2.yaml} (shipped with Cantera):
\begin{lstlisting}[language={}]
units: {length: cm, time: s, quantity: mol, activation-energy: cal/mol}

phases:
- name: ohmech
  thermo: ideal-gas
  elements: [O, H, Ar, N]
  species: [H2, H, O, O2, OH, H2O, HO2, H2O2, AR, N2]
  kinetics: gas
  transport: mixture-averaged

species:
- name: OH
  composition: {O: 1, H: 1}
  thermo:
    model: NASA7
    temperature-ranges: [200.0, 1000.0, 3500.0]
    data:                     # a1 ... a7, low range then high range
    - [3.99201543, -2.40131752e-03, 4.61793841e-06, -3.88113333e-09, 1.3641147e-12,
      3615.08056, -0.103925458]
    - [3.09288767, 5.48429716e-04, ...]
  transport: {model: gas, geometry: linear, well-depth: 80.0, diameter: 2.75}

reactions:
- equation: H + O2 <=> O + OH          # chain branching R2
  rate-constant: {A: 2.65e+16, b: -0.6707, Ea: 1.7041e+04}
- equation: H + O2 + M <=> HO2 + M     # termination R5
  type: three-body
  rate-constant: {A: 2.8e+18, b: -0.86, Ea: 0.0}
  efficiencies: {O2: 0.0, H2O: 0.0, N2: 0.0, AR: 0.0}
- equation: 2 OH (+M) <=> H2O2 (+M)
  type: falloff
  low-P-rate-constant: {A: 2.3e+18, b: -0.9, Ea: -1700.0}
  high-P-rate-constant: {A: 7.4e+13, b: -0.37, Ea: 0.0}
  Troe: {A: 0.7346, T3: 94.0, T1: 1756.0, T2: 5182.0}
\end{lstlisting}
\begin{itemize}
  \item \code{units}: everything below is in cm, mol, cal; Cantera converts to SI.
  \item NASA7 \code{data}: the $a_1\dots a_7$ of \secref{sec:nasa}, for $200$--$1000$ and $1000$--$\SI{3500}{\kelvin}$.
        $a_6=3615$ encodes OH's positive heat of formation.
  \item \code{transport}: Lennard-Jones well depth $\epsilon/k_B$ [K] and diameter $\sigma$ [\AA], used
        by kinetic theory (\chapref{ch:transport}).
  \item R2's $E_a=\SI{17041}{cal/mol}=\SI{71.3}{\kilo\joule\per\mole}$: the steep line of \figref{fig:arrh}.
  \item In R5 the efficiencies of \ce{O_2}, \ce{H_2O}, \ce{N_2}, Ar are zero because those colliders have
        their own explicit reactions (\code{H + O2 + N2 <=> HO2 + N2}, \dots). Do not count them twice.
\end{itemize}

\section{Equilibrium}
\label{sec:ct-equil}
\begin{lstlisting}
gas.set_equivalence_ratio(phi, "H2", "O2:1, N2:3.76")
gas.TP = T_u, p
gas.equilibrate("HP")          # adiabatic, constant pressure
T_ad = gas.T
\end{lstlisting}
\begin{keypoint}
Compute the equilibrium state \emph{before} any kinetic simulation. It is the value every reactor and
flame must approach, and the sanity check for all of them.
\end{keypoint}

\section{\texttt{Solution} loops vs.\ \texttt{SolutionArray}}
\label{sec:ct-solarray}

When states are independent (parameter sweeps, post-processing), a \code{SolutionArray} moves the loop
into C++:
\begin{lstlisting}
import numpy as np
phi = np.linspace(0.4, 3.0, 200)
states = ct.SolutionArray(gas, shape=phi.shape)
states.TP = 300.0, ct.one_atm
states.set_equivalence_ratio(phi, "H2", "O2:1, N2:3.76")
states.equilibrate("HP")
T_ad = states.T                          # numpy array, same shape as phi
\end{lstlisting}
Use a plain loop when each step depends on the previous one: time integration, or flame sweeps that
reuse the previous solution. For 2D maps (Part~1, Task~4) use a 2D shape, \code{SolutionArray(gas, shape=(n1, n2))}, and
\code{np.meshgrid} for the inputs.

\section{Reactors}
See \chapref{ch:reactor}. \code{IdealGasConstPressureReactor} (flames, open systems) or
\code{IdealGasReactor} (closed bomb), wrapped in a \code{ReactorNet} and advanced with
\code{net.step()} or \code{net.advance(t)}.

\section{1D flames}

\begin{table}[H]
\centering
\begin{tabular}{lll}
\toprule
class & physics & project \\
\midrule
\code{FreeFlame} & freely propagating premixed flame, $\SL$ eigenvalue & Part 3 \\
\code{CounterflowDiffusionFlame} & opposed fuel/oxidiser jets & Part 2 \\
\code{BurnerFlame} & burner-stabilised premixed (heat loss to burner) & -- \\
\code{CounterflowPremixedFlame} & opposed premixed jets (stretch) & -- \\
\bottomrule
\end{tabular}
\end{table}

Useful attributes after \code{solve}: \code{grid} [m], \code{T}, \code{velocity}, \code{X}, \code{Y}
($K\times N$ arrays), \code{heat\_release\_rate}, \code{density}, \code{cp}, \code{thermal\_conductivity},
\code{mixture\_fraction("Bilger")} (counterflow), \code{strain\_rate("max")}, \code{extinct()}.

\section{Grid refinement: the most important setting}
\label{sec:ct-refine}

Cantera starts from a coarse grid and adds points where the solution is under-resolved, according to
three criteria applied to every solution component $\varphi$:

\begin{figure}[H]
\centering
\begin{tikzpicture}[xscale=1.1]
  \begin{scope}
    \draw[->] (0,0) -- (3.5,0); \foreach \x in {0.3,0.6,0.9,2.2,3.2} \fill (\x,0) circle (1.6pt);
    \draw[<->,warn] (0.9,-0.3) -- (2.2,-0.3) node[midway,below,lbl] {too large vs.\ neighbour};
    \node[lbl,align=center] at (1.75,1.3) {\textbf{ratio}\\$\Delta x_{j+1}/\Delta x_j\le$ ratio};
  \end{scope}
  \begin{scope}[xshift=4.6cm]
    \draw[->] (0,0) -- (3.5,0);
    \draw[thick,hot] (0.2,0.1) -- (1.4,0.15) -- (2.0,0.95) -- (3.3,1.0);
    \foreach \x/\y in {0.2/0.1,1.4/0.15,2.0/0.95,3.3/1.0} \fill (\x,\y) circle (1.6pt);
    \draw[<->,warn] (2.3,0.15) -- (2.3,0.95) node[midway,right,lbl] {$\Delta\varphi$};
    \node[lbl,align=center] at (1.75,1.6) {\textbf{slope}\\$|\Delta\varphi|\le$ slope$\cdot$range($\varphi$)};
  \end{scope}
  \begin{scope}[xshift=9.2cm]
    \draw[->] (0,0) -- (3.5,0);
    \draw[thick,hot] (0.2,0.1) -- (1.4,0.15) -- (2.4,1.0) -- (3.3,1.05);
    \foreach \x/\y in {0.2/0.1,1.4/0.15,2.4/1.0,3.3/1.05} \fill (\x,\y) circle (1.6pt);
    \draw[warn,thick] (1.4,0.15) circle (6pt); \draw[warn,thick] (2.4,1.0) circle (6pt);
    \node[lbl,align=center] at (1.75,1.6) {\textbf{curve}\\$|\Delta(\dd\varphi/\dd x)|\le$ curve$\cdot$range};
  \end{scope}
\end{tikzpicture}
\caption{Refinement criteria. Smaller \code{slope} and \code{curve} mean more points in steep and
strongly curved regions. \code{prune} removes points where they are no longer needed.}
\end{figure}

\begin{lstlisting}
flame.set_refine_criteria(ratio=2.0, slope=0.005, curve=0.01, prune=0.0)  # research grade
flame.set_max_grid_points(flame.flame, 10000)
\end{lstlisting}

\begin{warning}[the defaults are not research grade]
The ASEC tutorial page shows a lean \ce{H_2} flame ($\phi=0.5$) where the default refinement gave only 27 points,
with errors of $\sim9\%$ in $T_b$, $\sim11\%$ in $Y_{\ce{H_2O}}$ and $\sim25\%$ in $\SL$, while the solver
reported success. \figref{fig:grid} shows the same lesson for lean methane. Always:
\begin{enumerate}
  \item print \code{len(flame.grid)};
  \item check the end state against \code{equilibrate("HP")} ($<1\%$);
  \item tighten \code{slope}/\code{curve} until $\SL$ changes by less than your tolerance;
  \item \textbf{lighter fuels need tighter criteria}: hydrogen flames are thin in the reaction zone yet
        have long diffusive tails.
\end{enumerate}
Do not use 1D flames for NO: thermal NO forms over milliseconds in the post-flame gas, far beyond the
domain. Use a reactor or plug-flow model downstream.
\end{warning}

A practical strategy is to solve with loose criteria first, then tighten and re-solve: the coarse
solution is an excellent initial guess (\secref{sec:newton}).
\begin{lstlisting}
for slope, curve in [(0.1, 0.2), (0.05, 0.1), (0.02, 0.04), (0.01, 0.02)]:
    flame.set_refine_criteria(ratio=2, slope=slope, curve=curve)
    flame.solve(loglevel=0, refine_grid=True, auto=False)
    print(len(flame.grid), flame.velocity[0])
\end{lstlisting}

\section{Sweeps, continuation, saving}
\label{sec:ct-sweeps}

\begin{lstlisting}
flame.solve(loglevel=0, auto=True)               # first point: robust
results = []
for phi in phi_values:                            # neighbours: start from previous solution
    gas.set_equivalence_ratio(phi, "H2", "O2:1, N2:3.76")
    gas.TP = T_u, p
    flame.inlet.T = T_u
    flame.inlet.X = gas.X
    flame.solve(loglevel=0, auto=False)
    results.append((phi, flame.velocity[0]))
    flame.save("flames.yaml", name=f"phi_{phi:.2f}", overwrite=True)

flame.restore("flames.yaml", name="phi_1.00")    # reload later without re-solving
\end{lstlisting}

\begin{itemize}
  \item Sweep \emph{outwards} from the easy case ($\phi\approx1$) towards the limits; the solver
        struggles near flammability limits.
  \item If a solve fails, catch \code{ct.CanteraError}, go back to the last good solution and take a
        smaller parameter step (this is how \figref{fig:ext} was made).
  \item Change the transport model on a converged solution and re-solve. This is faster than starting
        over: \code{flame.transport\_model = "unity-Lewis-number"; flame.solve(auto=False)}.
\end{itemize}

\section{API notes for Cantera 3.2}
\begin{itemize}
  \item Transport models: \code{"mixture-averaged"}, \code{"multicomponent"}, \code{"unity-Lewis-number"}.
  \item Reactors: pass \code{clone=True/False} explicitly; use \code{reactor.phase} (not \code{.thermo}).
  \item Flames: \code{flame.save}/\code{restore} write YAML (or HDF5 with \code{.h5});
        \code{flame.to\_pandas()} gives a DataFrame for plotting.
  \item Counterflow component names: \code{velocity}, \code{spreadRate}, \code{T}, \code{Lambda}, species.
\end{itemize}
